\renewcommand{\arraystretch}{1}

\fancyhead[L]{ARSE32}
\fancyhead[R]{1 NSI}
\titre{}{ARSE32 - TP LINUX}


\section{A la découverte de Raspbian}

\subsection{Présentation}

\begin{minipage}{9cm}
La barre des tâches comporte une partie gauche (ci-dessous), une zone centrale affichant les fenêtres
ouvertes sur le bureau ou mises en réduction et une partie droite donnant des informations sur les
connexions réseau actives ou encore l'heure.
\end{minipage}
\begin{minipage}{9cm}
\begin{center}
\includegraphics[scale=0.25]{rasp2.ps}
\end{center}
\end{minipage}




\subsection{Présentation du Terminal}

\exer 

\begin{enumerate}[A.]
	\item Ouvrez une fenêtre "LX Terminal"

\begin{center}
\includegraphics[scale=0.4]{exo1.ps}
\end{center}



Le système de fichiers de Raspbian (Linux) est organisé à partir d'un point de départ appelé root, racine ou encore "/". Sous cette racine se déploient des répertoires contenant les fichiers et programmes nécessaires au système d’exploitation.\\


	\item Taper la commande \fbox{ pwd }. Que permet-elle de connaître? \dotfill
	
	\dotfill
	
	\item Remonter au plus haut niveau de l'arborescence en tapant la commande \fbox{ cd / }. (cd veut dire "change working directory").
	
	\item Taper la commande \fbox{ ls }. Que permet-elle de connaître? \dotfill
	
	\dotfill
	
	\item Dans ce répertoire, il y a des fichiers et des répertoires. Comment les reconnaître? \dotfill
	
	\dotfill
	
\begin{center}
\includegraphics[scale=0.5]{exo1E.ps}
\end{center}	

	\item Vous pouvez visualiser l'arborescence en tapant \fbox{ tree -L 1 }. On visualise alors l'arborescence de niveau 1. Essayer en remplaçant "1" par "2". Qu'observe-t-on? \dotfill
	
	\dotfill
	
	\item Déplacez-vous dans l'arborescence avec la commande \fbox{ cd } suivi du répertoire dans lequel vous voulez vous déplacer. Allez dans le répertoire "home". Comment savez-vous que vous êtes bien dans le répertoire "home"? \dotfill
	
	\dotfill
	
	\item Quels sont les répertoires présents dans le répertoire "home"? \dotfill
	
\end{enumerate}

\medskip

\cadre{
\prop Le système de fichiers racine (root file system), soit le système de fichiers primaire est associé au répertoire le plus haut / : 

\begin{minipage}{10cm}
\begin{center}
\includegraphics[scale=0.45]{arbo.ps}
\end{center}
\end{minipage}
\begin{minipage}{8cm}
\begin{liste}
	\item \fbox{ . } est le répertoire courant
	\item \fbox{ .. } est le répertoire parent
	\item \fbox{ \textasciitilde } est le répertoire "home" de l'utilisateur
	\item \fbox{ / } est la racine du système.
\end{liste}
\end{minipage}}


\medskip

\exer

\begin{enumerate}[A.]
	\item Positionnez-vous dans le répertoire \verb+\etc+
	\item Taper la commande \fbox{ ls - al }. Que permet-elle de visualiser? \dotfill
	
	\dotfill
	
	\item Comment reconnait-on les répertoires des fichiers autrement que par la couleur de l'écriture? \dotfill
	
\dotfill

	\item Taper la commande \fbox{ ls - al | more } afin de pouvoir visualiser le contenu du répertoire. Enfoncez la barre d'espace pour activer le défilement.
	\item D'après-vous, \fbox{ python }  est-il un fichier ou un répertoire? Comment le vérifier? \dotfill
	
	\dotfill
	
	\item Déplacez-vous dans le répertoire "xml".
	
	\item Quelle commande doit-on taper pour se déplacer vers \fbox{ /etc/vim } 
		\begin{enumerate}[1.]
			\item en chemin relatif: \dotfill
			\item en chemin absolu: \dotfill
	\end{enumerate}		 
\end{enumerate}


\subsection{Modifier l'arborescence}

\subsubsection{Créer un répertoire}

La commande \fbox{ mkdir } (make directory) permet de créer un répertoire dans le répertoire courant. La commande est \fbox{ mkdir nom\_du\_repertoire }

\medskip

\exer

\begin{enumerate}[A.]
	\item Déplacez-vous dans le répertoire \fbox{ /home/pi }
	\item Notez ici le contenu de ce répertoire. \dotfill
	
	\dotfill
	
	\item Créer un répertoire \fbox{ Mes\_doc\_NSI } en utilisant la commande \fbox{ mkdir }. Commande tapée: \dotfill
	
	\dotfill
	
	\item Déplacer vous dans le répertoire \fbox{ Mes\_doc\_NSI }. Commande tapée: \dotfill
	
	\item Créez les répertoire suivants: \fbox{ Mes\_ photos }, \fbox{ Mes\_doc\_pdf }, \fbox{ Mes\_sites\_web }
	\item Taper la commande \fbox{mkdir -p Mes\_algorithmes/programmes }. Quel est son effet? \dotfill
	
	\dotfill
	
	\dotfill
\end{enumerate}

\bigskip

\subsubsection{Créer un fichier avec un éditeur de texte}

\exer

\begin{enumerate}[A.]
	\item Déplacez-vous dans le répertoire  \fbox{ /home/pi/ Mes\_doc\_NSI/ Mes\_algorithmes/programmes } en utilisant un chemin relatif
	\item Tapez la commande \fbox{ nano compteur.py}. Vous parvenez alors dans la fenêtre de l'éditeur de texte « nano ». nano est un éditeur de texte tout simple. Voici les commandes essentielles :
	
	\begin{center}
\includegraphics[scale=0.6]{exo6A.ps}
\end{center}
\item Écrire un programme en langage Python qui initialise un compteur qui compte et affiche les entiers de 0 à 25 inclus. 
\item Sauvegarder le fichier: CTRL+O puis ENTREE

	\begin{center}
\includegraphics[scale=0.6]{exo6D.ps}
\end{center}

\item Quitter "nano": CTRL+X
\item Vérifier que votre fichier a été créé dans le répertoire "programme"
\item Visualiser le contenu du fichier en tapant \fbox{ more compteur.py }. Le conteu s'affiche en lecture seule.
\item Exécuter votre programme en ligne de commande grâce à: \fbox{ python3 compteur.py }.
\end{enumerate}
	

\subsubsection{Copier des fichiers}


La copie de fichiers ou de répertoires se fait grâce à la commande \fbox{ cp }. La syntaxe est la suivante: \fbox{ cp [OPTION] <fichier> <destination>}

\medskip

Vous pouvez accéder aux options de la commande avec: \fbox{ cp - -help }.

		\begin{center}
\includegraphics[scale=0.45]{cp1.ps}
\end{center}

\exer

\begin{enumerate}[A.]
	\item Créer le répertoire \fbox{ /home/pi/ Mes\_doc\_NSI/ Mes\_algorithmes/programmes\_python }
	\item Copier le fichier \fbox{ compteur.py } dans le répertoire \fbox{ programmes\_python/} et le renommer \fbox{ compteur2.py }.
	\item Quelle commande avez-vous tapé? \dotfill
	\item Modifier le programme \fbox{ compteur2.py} afin qu'il compte de 0 à 13 puis de 13 à 0.
\end{enumerate}

\medskip

	
\cadre{\prop Copier un fichier vers un autre fichier: 

Si le fichier n'existe pas, le fichier est copié. Si la destination existe, le fichier est écrasé.}
	
	\newpage
	
	\exer
	
	\begin{enumerate}[A.]
		\item Positionnez-vous dans le répertoire \fbox{ /Mes\_doc\_NSI/Mes\_algorithmes/programmes } et créez un nouveau fichier nommé \fbox{ remplissLilste.py } dans lequel vous taperez un programme permettant de remplir une liste contenant les 10 premières puissances de 2.
		\item Copiez les deux fichiers \fbox{ compteur.py } et \fbox{ remplissListe.py } dans 
		
		\fbox{ /Mes\_doc\_NSI/Mes\_algorithmes/programmes\_python } avec la commande suivante: 
		
		\fbox{ cp -iu fichier1\_1 fichier\_2 ... fichier\_n repertoire }.
		
	\end{enumerate}
	
	
	
	\subsubsection{Copier un répertoire}
	
	La syntaxe est la suivante: \fbox{ cp [options] <repertoire\_source> <repertoire\_destination>}
	
	\begin{center}
\includegraphics[scale=0.5]{cp2.ps}
\end{center}

\medskip

\exer

\begin{minipage}{12cm}
\begin{enumerate}[A.]
	\item Au sein du répertoire \fbox{ /home/pi/ Mes\_doc\_NSI }, créez le répertoire \fbox{ /Mes\_programmes }
	\item Copiez le répertoire \fbox{ /Mes\_algorithmes } ainsi que son arborescence vers le répertoire 
	
	\fbox{ /home/pi/ Mes\_doc\_NSI/Mes\_programmes }
	\item Vérifiez la bonne exécution de l’opération en utilisant une seule commande.
	\item Vérifier que votre arborescence depuis \fbox{ Mes\_doc\_NSI} est la suivante:
\end{enumerate}	
\end{minipage}
\begin{minipage}{6cm}
\begin{center}
\includegraphics[scale=0.65]{exo9D.ps}
\end{center}
\end{minipage}

	
\subsubsection{Déplacer et renommer un fichier}

La syntaxe est la suivante: \fbox{ mv [options] <source> <destination> }

		\begin{center}
\includegraphics[scale=0.75]{deplacer.ps}
\end{center}

\medskip

\exer

\begin{enumerate}[A.]
	\item Positionnez-vous dans le répertoire \fbox{ /home/pi/ Mes\_doc\_NSI/Mes\_programmes }
	\item Déplacez le répertoire 
	\fbox{ Mes\_algorithmes/programmes\_python } vers 
	\fbox{ Mes\_programmes } en le renommant 
	\fbox{ Mes\_programmes/Mes\_programmes\_python }
	
	\item Vérifiez la bonne exécution de l'opération.
	
	\item Positionnez-vous dans \fbox{  /home/pi/ Mes\_doc\_NSI/Mes\_programmes/Mes\_programmes\_python }
	
	\item Renommez 
	\fbox{ compteur.py } en 
	\fbox{   compteur0\_25.py } et 
	\fbox{ compteur2.py } en 
	\fbox{   compteur0\_13\_0.py }
	
	\item Depuis le répertoire \fbox{ /home/pi/ Mes\_doc\_NSI/Mes\_programmes/Mes\_programmes\_python }, déplacez le fichier \fbox{  remplissListe.py } vers le répertoire \fbox{  /home/pi/ Mes\_doc\_NSI/Mes\_doc\_pdf } en utilisant un chemin relatif
	
	\item Vérifier que votre arborescence est la suivante:
	
			\begin{center}
\includegraphics[scale=0.8]{exo10G.ps}
\end{center}
\end{enumerate}


\subsubsection{Supprimer des fichiers}

La syntaxe est la suivante: \fbox{ rm [options] <nom\_de\_fichier> }

		\begin{center}
\includegraphics[scale=0.8]{suppr.ps}
\end{center}

\medskip

\exer

\begin{enumerate}[A.]
	\item Supprimez le répertoire \fbox{ Mes\_doc\_NSI/ Mes\_algorithmes } avec son contenu
	\item Supprimez le répertoire \fbox{  Mes\_doc\_NSI/Mes\_programmes/Mes\_algorithmes } avec son contenu
	\item Supprimez le répertoire \fbox{ Mes\_doc\_NSI/Mes\_programmes/Mes\_algorithmes/programmes }
avec son contenu
	\item Vérifiez le contenu de votre arborescence. Celle-ci doit être similaire à :
	
			\begin{center}
\includegraphics[scale=0.8]{exo11D.ps}
\end{center}
\end{enumerate}

\subsubsection{Afficher le contenu d'un fichier}

La syntaxe est la suivante: \fbox{ cat [options] <nom\_de\_fichier\_1> <nom\_de\_fichier\_2>}

		\begin{center}
\includegraphics[scale=0.8]{affichagecontenu.ps}
\end{center}

\medskip

\exer Afficher le contenu des deux programmes \fbox{ compteur0\_25.py } et \fbox{ compteur0\_13\_0.py }.


\section{Gestion des utilisateurs et des groupes}

Il est possible d'utiliser la commande "ls" avec l'option "-l" afin d'avoir des informations supplémentaires. 

		\begin{center}
\includegraphics[scale=0.7]{gestion1.ps}
\end{center}


\subsection{Changer les droits utilisateurs}

Il est important de ne pas perdre de vu que l'utilisateur "root" a la possibilité de modifier les droits de tous les utilisateurs.\\

Le propriétaire d'un fichier peut modifier les permissions d'un fichier ou d'un répertoire à l'aide de la commande \fbox{ chmod }. Pour utiliser cette commande, il est nécessaire de connaitre certains symboles :

\begin{liste}
	\item les symboles liés aux utilisateurs : "u" correspond au propriétaire, "g" correspond au groupe lié au fichier (ou au répertoire), "o" correspond aux autres utilisateurs et "a" correspond à "tout le monde" (permet de modifier "u", "g" et "o" en même temps)
    \item les symboles liés à l'ajout ou la suppression des permissions : "+" on ajoute une permission, "-" on supprime une permission, "=" les permissions sont réinitialisées (permissions par défaut)
    \item les symboles liés aux permissions : "r" : lecture, "w" : écriture, "x" : exécution.
\end{liste}

\bigskip

La syntaxe est donc la suivante: \fbox{ chmod [ u g o a ] [ + - = ] [ r w x ] nom\_du\_fichier }

\bigskip

\exemple Par exemple la commande \fbox{ chmod o+w toto.txt } attribuera la permission "écriture" pour le fichier "toto.txt" "aux autres utilisateurs" 

\newpage


\exer 

\begin{enumerate}[A.]
	\item Créer un nouveau fichier dans le répertoire \fbox{ Mes\_docs\_NSI } nommé \fbox{ toto.txt } avec la commande \fbox{ touch toto.txt }.
	\item Vérifier les permissions du fichier créé.
		\begin{liste}
			\item Propriétaire:\dotfill
			\item Groupe:\dotfill
			\item Autres:\dotfill
		\end{liste}
	\item Autoriser les autres à exécuter ce fichier. Commande: \dotfill
	\item Empêcher le groupe de lire ce fichier. Commande: \dotfill
	\item Autoriser tout le monde à écrire dans ce fichier. Commande: \dotfill
	\item Empêcher le propriétaire d'exécuter le fichier. Commande: \dotfill
	
	
\end{enumerate}

\section{Gestion des utilisateurs}

\subsection{Ajouter un utilisateur}

% passwd root
% sudo passwd root

Les systèmes de type "UNIX" sont des systèmes multi-utilisateurs, plusieurs utilisateurs peuvent donc partager un même ordinateur, chaque utilisateur possédant un environnement de travail qui lui est propre.\\

Dans l'arborescence actuelle il y a u utilisateur nommé \fbox{ pi }\\

Il est possible d'ajouter un utilisateur.

\begin{enumerate}
	\item Changer le mot de passe root: \fbox{ sudo passwd root } et taper un nouveau mot de passe root. Pour que tout le monde ait le même, taper \fbox{ root }
	\item Accéder à la console en mode super-utilisateur: \fbox{ su - }
	\item Ajouter un utilisateur \fbox{ adduser nom\_utilisateur }. Ajouter un utilisateur nommé \fbox{ toto }, mot de passe \fbox{ toto }
	\item Vérifier que le nouvel utilisateur a bien été créé en allant dans \fbox{ home }
	\item Déconnectez vous et logger vous sous le nom \fbox{ toto }
\end{enumerate}

\subsection{Ajouter un utilisateur à un groupe}

\begin{enumerate}
\item Identifiez vous avec le compte \fbox{ pi } mot de passe \fbox{ raspberry }.
	\item Accéder à la console en mode super-utilisateur: \fbox{ su - }
	\item créer un groupe nommé \fbox{ eleve } avec la commande \fbox{ sudo addgroup eleve }
	\item Ajouter toto au groupe \fbox{ eleve } avec la commande \fbox{ usermod -aG toto eleve }
	\end{enumerate}
	
\subsection{Suppression d'un utilisateur}

\begin{enumerate}
	\item Accéder à la console en mode super-utilisateur: \fbox{ su - }
	\item supprimer toto du groupe \fbox{ eleve } avec la commande \fbox{ sudo deluser toto }
	\end{enumerate}
	
\subsection{Suppression d'un groupe}

\begin{enumerate}
	\item Accéder à la console en mode super-utilisateur: \fbox{ su - }
	\item supprimer le groupe \fbox{ eleve }  avec la commande \fbox{ sudo delgroup eleve }
	\end{enumerate}
	
\import{../../../../preambule}{raz_compteurs.tex}
\newpage